\chapter{Overview}
\section{Purpose}
This guide explains how EFC approaches the design, implementation, testing, commissioning, documentation, operation, and maintenance of industrial control systems. It is written primarily for engineers who are new to EFC projects and is intended to help them build systems that are safe, supportable, repeatable, secure, and consistent.

\section{Scope}
The guide covers PLCs, remote I/O, industrial control panels, operator workstations, HMI/SCADA systems, industrial networks, instrumentation interfaces, packaged-equipment integration, and the related project documents. It applies to internal engineering work and to contractors unless a project specification gives a more stringent requirement.

This document does not replace the adopted electrical code, the requirements of the Authority Having Jurisdiction (AHJ), a project risk assessment, a process hazard analysis, or a machine-safety/functional-safety assessment. The project engineer must identify the governing jurisdiction and applicable editions at project initiation.

\section{How to Use This Guide}
\begin{itemize}[itemsep=2pt]
  \item \textbf{Must} identifies a requirement that is necessary for safety, compliance, security, or EFC's basic engineering practices.
  \item \textbf{Should} identifies the normal approach. If a project needs something different, document the reason.
  \item \textbf{May} indicates a permitted option.
\end{itemize}

\section{Governance}
The Controls Engineering function owns this standard. The Controls Engineering Manager or delegated technical authority approves revisions and project deviations. The standard must be reviewed at least annually and whenever a referenced code, core platform, cybersecurity requirement, or approved hardware family changes materially.

Project deviations must be recorded in the design basis and include the affected requirement, technical justification, risk evaluation, compensating measures, approver, and expiration or closeout condition. A deviation does not modify this standard for other projects.

\section{Controlled Layers}
Controls documentation must be maintained in four layers:
\begin{enumerate}[itemsep=2pt]
  \item Company-wide mandatory engineering requirements.
  \item Approved hardware and software catalog, including lifecycle status.
  \item Version-controlled reusable PLC, SCADA, and drawing libraries.
  \item Project-specific design, test, configuration, and validation records.
\end{enumerate}
